\section{Limits}
\label{sec:limits}

A derivative describes change at a single point, but it is constructed from changes measured over a finite interval. The mathematical bridge between these two descriptions is the \emph{limit}. A limit allows us to study the value approached by a function even when the function is not defined at the point being approached.

\subsection{Approaching a Point}

Consider the function
\[
f(x)=\frac{x^2-1}{x-1}.
\]
At \(x=1\), the denominator vanishes, so the expression is undefined. Nevertheless, its behavior near \(x=1\) is regular. Since
\[
x^2-1=(x-1)(x+1),
\]
we may simplify the expression for every \(x\neq1\):
\[
f(x)=x+1,
\qquad x\neq1.
\]
As \(x\) approaches \(1\), the values of \(f(x)\) approach \(2\). We write
\[
\boxed{
\lim_{x\to1}\frac{x^2-1}{x-1}=2.
}
\]

\begin{bookdefinition}[Informal Meaning of a Limit]
The statement
\[
\lim_{x\to a}f(x)=L
\]
means that the values of \(f(x)\) can be made arbitrarily close to \(L\) by choosing \(x\) sufficiently close to \(a\), while allowing \(x\neq a\).
\end{bookdefinition}

The distinction between a function value and a limiting value is fundamental. In the example above, \(f(1)\) is undefined, yet the limit exists and equals \(2\). A limit depends on the behavior of the function \emph{near} the point, not necessarily on its value \emph{at} the point.

\begin{bookinsight}[A Punctured Neighborhood]
When evaluating \(\lim_{x\to a}f(x)\), one examines values of \(x\) in a small neighborhood around \(a\) while excluding the point \(x=a\) itself. For this reason, limits can describe removable holes, discontinuities, and singular expressions without requiring the original formula to be defined at the limiting point.
\end{bookinsight}

\subsection{The Precise Definition}

The intuitive phrase ``arbitrarily close'' is made rigorous using two positive numbers. The number \(\varepsilon\) measures the permitted error in the function value, while \(\delta\) measures how close the input must remain to the point \(a\).

\begin{bookdefinition}[The \(\varepsilon\)--\(\delta\) Definition]
Let \(f\) be defined on a punctured neighborhood of \(a\). We say that
\[
\lim_{x\to a}f(x)=L
\]
if, for every \(\varepsilon>0\), there exists a \(\delta>0\) such that
\[
0<|x-a|<\delta
\quad\Longrightarrow\quad
|f(x)-L|<\varepsilon.
\]
\end{bookdefinition}

The logical order is important. An arbitrary accuracy \(\varepsilon\) is specified first. We must then find a corresponding input tolerance \(\delta\) that guarantees the required accuracy. The value of \(\delta\) may depend on \(\varepsilon\), but it may not depend on the particular value of \(x\).

\begin{bookexample}[Proving a Linear Limit]
Prove that
\[
\lim_{x\to2}(3x-1)=5.
\]
We begin with the desired condition
\[
|(3x-1)-5|<\varepsilon.
\]
Since
\[
|(3x-1)-5|=|3x-6|=3|x-2|,
\]
it is sufficient to require
\[
|x-2|<\frac{\varepsilon}{3}.
\]
Therefore choose
\[
\delta=\frac{\varepsilon}{3}.
\]
Then
\[
0<|x-2|<\delta
\]
implies
\[
|(3x-1)-5|=3|x-2|<3\delta=\varepsilon.
\]
Hence, by the definition of a limit,
\[
\boxed{\lim_{x\to2}(3x-1)=5.}
\]
\end{bookexample}

\subsection{One-Sided Limits}

A function may approach different values from the left and from the right. These behaviors are described by one-sided limits:
\[
\lim_{x\to a^-}f(x)
\qquad\text{and}\qquad
\lim_{x\to a^+}f(x).
\]
The superscript \(-\) indicates that \(x<a\), while the superscript \(+\) indicates that \(x>a\).

\begin{bookprinciple}[Existence of a Two-Sided Limit]
The two-sided limit
\[
\lim_{x\to a}f(x)=L
\]
exists if and only if both one-sided limits exist and are equal:
\[
\lim_{x\to a^-}f(x)
=
\lim_{x\to a^+}f(x)
=L.
\]
\end{bookprinciple}

\begin{bookexample}[A Jump Discontinuity]
Let
\[
f(x)=
\begin{cases}
-1, & x<0,\\
 1, & x\geq0.
\end{cases}
\]
Then
\[
\lim_{x\to0^-}f(x)=-1,
\qquad
\lim_{x\to0^+}f(x)=1.
\]
Because the one-sided limits are unequal, the two-sided limit \(\lim_{x\to0}f(x)\) does not exist.
\end{bookexample}

\subsection{Algebra of Limits}

Once basic limits are known, more complicated limits may be evaluated using limit laws.

\begin{bookproperty}[Limit Laws]
Suppose
\[
\lim_{x\to a}f(x)=L,
\qquad
\lim_{x\to a}g(x)=M,
\]
and let \(c\) be a constant. Then
\[
\lim_{x\to a}\bigl(f(x)+g(x)\bigr)=L+M,
\]
\[
\lim_{x\to a}\bigl(cf(x)\bigr)=cL,
\]
\[
\lim_{x\to a}\bigl(f(x)g(x)\bigr)=LM,
\]
and, provided \(M\neq0\),
\[
\lim_{x\to a}\frac{f(x)}{g(x)}=\frac{L}{M}.
\]
For positive integers \(n\),
\[
\lim_{x\to a}[f(x)]^n=L^n.
\]
\end{bookproperty}

These rules justify direct substitution for polynomials and for rational functions whose denominators do not vanish at the limiting point.

\begin{bookexample}[A Limit Evaluated by Algebra]
Evaluate
\[
\lim_{x\to3}\frac{x^2+2x-1}{x+1}.
\]
Because the denominator is nonzero at \(x=3\), the limit laws permit direct substitution:
\[
\lim_{x\to3}\frac{x^2+2x-1}{x+1}
=
\frac{3^2+2(3)-1}{3+1}
=
\frac{14}{4}
=
\frac{7}{2}.
\]
\end{bookexample}

\subsection{The Squeeze Principle}

Some limits are difficult to evaluate directly but can be trapped between two simpler functions.

\begin{bookprinciple}[Squeeze Principle]
Suppose that, for all \(x\) sufficiently close to \(a\),
\[
g(x)\leq f(x)\leq h(x),
\]
and suppose that
\[
\lim_{x\to a}g(x)
=
\lim_{x\to a}h(x)
=L.
\]
Then
\[
\lim_{x\to a}f(x)=L.
\]
\end{bookprinciple}

A central example is
\[
\boxed{
\lim_{x\to0}\frac{\sin x}{x}=1,
}
\]
where \(x\) is measured in radians. This result is one of the basic limits used to derive the derivatives of trigonometric functions.

\subsection{Limits and Continuity}

A function is continuous at a point when its limiting value agrees with its actual value.

\begin{bookdefinition}[Continuity at a Point]
A function \(f\) is continuous at \(x=a\) if
\[
\lim_{x\to a}f(x)=f(a).
\]
Equivalently, all three of the following conditions must hold:
\begin{enumerate}
\item \(f(a)\) is defined;
\item \(\lim_{x\to a}f(x)\) exists;
\item the limit equals the function value.
\end{enumerate}
\end{bookdefinition}

Continuity expresses the absence of a break, jump, or hole at the point under consideration. Most elementary physical models assume that position, temperature, density, and field variables vary continuously over the scales at which the model is applied.

\begin{bookremark}[Continuity Is a Modeling Assumption]
Continuity is not guaranteed by nature at every scale. A continuum model may treat matter as smoothly distributed even though matter is atomistic microscopically. The usefulness of continuity therefore depends on the physical scale and the approximation being made.
\end{bookremark}

\subsection{Physical Meaning of the Limiting Process}

Suppose a particle moves along a line with position \(x(t)\). Its average velocity over a finite interval \(\Delta t\) is
\[
\bar v
=
\frac{x(t+\Delta t)-x(t)}{\Delta t}.
\]
To obtain the velocity at the instant \(t\), the interval must shrink toward zero:
\[
v(t)
=
\lim_{\Delta t\to0}
\frac{x(t+\Delta t)-x(t)}{\Delta t}.
\]
The interval never needs to be set equal to zero; doing so would produce division by zero. Instead, the ratio is examined for nonzero intervals that become arbitrarily small.

\begin{bookinsight}[Why Limits Matter]
An instantaneous rate of change cannot be measured directly from a zero-length interval. Limits retain the information contained in ratios over finite intervals while allowing those intervals to approach zero. This is the central mathematical idea behind differentiation.
\end{bookinsight}

With the limiting process now defined, we are prepared to construct the derivative and distinguish average change from instantaneous change.
